\chapter{Introduction} 
\label{chap:introduction} 

Industrial automation systems increasingly connect their operational technology to information technology networks. This convergence of Information Technology (IT) and Operational Technology (OT) links traditionally isolated control environments with corporate and external networks, and it turns the components at the boundary between the two domains into critical assets \autocite[2]{EHIE2020103166}. Industrial gateways occupy exactly this boundary. They translate and forward data between fieldbus protocols and higher-level services \autocite[277]{Gsell2026}, and they therefore expose the previously closed OT environment to threats that originate in the connected IT networks \autocite[199]{8586807}.

For a long time, these gateways were operated inside perimeter-protected zones, and their validation focused on functional behaviour, performance, and functional safety. Cybersecurity was not a mandatory part of this process. With Regulation (EU) 2023/1230, the European legislator changed this situation. The new Machinery Regulation introduces binding cybersecurity requirements and treats the protection against corruption and against malicious third parties as a prerequisite for the safety of machinery \autocite{EU2023_1230}. Cybersecurity thereby becomes a mandatory element of the conformity assessment, and the functional safety of a machine can no longer be evaluated independently of its cybersecurity \autocite{carl2026_maschinenverordnung}. The regulation is complemented by the harmonised draft standard prEN 50742, which provides the technical guidance for its implementation, and by the IEC 62443 series, which defines the component-level security requirements for industrial automation and control systems \autocite{prEN50742_2025, IEC62443-4-2}.

Although the regulatory expectations are now explicit, the established validation process does not yet meet them. Existing functional and safety validation verifies the intended behaviour of a device, but it does not systematically verify its security functions, and it does not provide traceable evidence that the implemented technical measures mitigate the relevant threats. This creates a gap between the regulatory expectations on the one hand and the currently established company methodology for functional testing on the other. The central problem addressed in this thesis is therefore the absence of a systematic, repeatable, and standards-based approach for the functional testing of security in industrial gateways at the physical IT/OT interface.

The task of this thesis is to close this gap for a representative device. The objective is to research security testing approaches, frameworks, and organisational practices, and to develop a systematic testing methodology that verifies the security functions of an industrial gateway, that maintains traceability between the regulatory requirements and the concrete test activities, and that can be reused for future product releases. The investigation focuses on the physical IT/OT interface and aims to provide empirical evidence of the effectiveness of the technical mitigation measures.
%TODO ist das wirklich mein Ziel
A constraint that shapes the whole work is that a security measure or a security test must never impair the functional safety of the device \autocite[17]{prEN50742_2025}.

The work is guided by three research questions. First, how can the cybersecurity obligations of Regulation (EU) 2023/1230 be translated into verifiable, component-level security tests for a safety-related industrial gateway? Second, which tests of an existing company security framework are applicable to an embedded gateway that exposes only a diagnostic interface, and how can non-applicable tests be identified before execution? Third, what can direct-access testing at the physical IT/OT interface verify about the effectiveness of the implemented security functions, and where are its boundaries?

The starting point of the work is the established functional test process of the industrial partner together with the newly introduced regulatory requirements. Methodologically, the thesis follows the Design Science Research Methodology, which structures the creation and the evaluation of an artifact that solves a practical engineering problem \autocite[54]{Peffers}. The artifact is a testing methodology that derives the required security tests in three stages. It extracts the security requirements from the Machinery Regulation, maps them through prEN 50742 to the component requirements of IEC 62443-4-2, and links each requirement to the affected assets and threats using threat-analysis and STRIDE principles \autocite[46]{doCarmo2025}. On this basis the methodology defines a two-level test catalogue, selects suitable testing techniques, and reuses an existing company test framework where it is applicable. The methodology is finally demonstrated on a real safety gateway in an isolated laboratory, and the results are assessed against the derived requirements and against the thesis objective. The target situation is a traceable testing methodology that facilitates the practical application of the regulatory requirements within an industrial manufacturing context and that supplies the evidence required for the technical documentation.

The remainder of this thesis is structured as follows. Chapter \ref{chap:fundamentals} establishes the technical, regulatory, and methodological foundations, including the relevant cybersecurity concepts, the IT/OT environment, the regulatory framework, and the established security testing approaches. Chapter \ref{chap:methodology} develops the testing methodology, from the requirement derivation and the traceability model to the selection of the testing approach, the design of the test environment, and the evaluation metrics. Chapter \ref{chap:implementation_results} applies the methodology to a physical Safe Remote I/O device, reports the collected evidence, verifies the derived requirements, and validates the methodology against the thesis objective. Chapter \ref{chap:summary_outlook} summarises the findings and gives an outlook on future work.


